% sections/01-intro.tex
\section{外侧度的引入}

\textbf{测度}, 作为一个如此抽象的名词, 但是其思想很简单, 就是去 "测量" 一个东西的大小, 例如长度、面积、体积等. 我们首先从外测度来入手.

\subsection{集合论前置}

在开始之前, 我们需要去了解一些基本的前置知识, 比如集合论的一些内容.

我们给定一个大集合 $X$.

\begin{definition}[上极限集]\label{def:suplim-set}
    设 $\{A_n\}_{n=1}^{+\infty}$ 是一个集合族, 则 $A_n$ 的上极限集定义为
    \[
    \limsup A_n = \bigcap_{n=1}^{+\infty} \bigcup_{k=n}^{+\infty} A_k
    \]
\end{definition}
\begin{definition}[下极限集]\label{def:inflim-set}
    设 $\{A_n\}_{n=1}^{+\infty}$ 是一个集合族, 则 $A_n$ 的下极限集定义为
    \[
    \liminf A_n = \bigcup_{n=1}^{+\infty} \bigcap_{k=n}^{+\infty} A_k
    \]
\end{definition}

用自然语言来说, 上极限集就是那些有无限多的 $n$ 使得 $x \in A_n$ 的点的集合, 而下极限集则是那些只有有限的的 $n$ 使得 $x \notin A_n$ 的点的集合.

换句话说, $\liminf A_n = \{x \in X :\exists N, \forall n \geq N, x \in A_n\}$, $\limsup A_n = \{x \in X :\forall N, \exists n \geq N, x \in A_n\}$.

不难看出, 我们总有下式成立:
\begin{equation}
    \liminf A_n \subseteq \limsup A_n \label{eq:liminf-limsup}
\end{equation}

当 \cref{eq:liminf-limsup} 中等号成立时, 我们称 $A_n$ 收敛到某个极限集 $A$.

而在集合的运算中, 我们经常用到集合的交, 并, 补, 差等操作, 所以自然的, 我们想在一个对这些运算封闭的系统里去研究这些集合. 我们称其为 $\sigma$-代数.

\begin{definition}[$\sigma$-代数]\label{def:sigma-algebra}

    $\mathscr{A} \subset \mathcal{P} = 2^X$ 称为 $\sigma$-代数, 如果满足:
    \begin{enumerate}
        \item $\emptyset \in \mathscr{A}, X \in \mathscr{A}$.
        \item 如果 $A \in \mathscr{A}$, 则 $A^c \in \mathscr{A}$.
        \item 如果 $\{A_n\}_{n=1}^{\infty} \subset \mathscr{A}$, 则 $\bigcup\limits_{n=1}^{\infty} A_n \in \mathscr{A}$.
    \end{enumerate}
\end{definition}

注意, 当他满足对差和并的封闭性时, 由 De-Morgan 定律可知, $\sigma$-代数也对交运算封闭.

对于 $\sigma$-代数, 可以举几个简单的例子:

\begin{example}
    平凡 $\sigma$-代数 $\{\emptyset, X\}$.
\end{example}
\begin{example}
    $2^{X}$ 是一个 $\sigma$-代数.
\end{example}
\begin{example}[一个有限集的例子]
    $X = \{a, b, c\}$, $A = \{\emptyset, \{a\}, \{b, c\}, X\}$ 是一个 $\sigma$-代数.
\end{example}
\begin{example}[Borel 集]
    包含 $\RR$ 中的所有开集的最小的$\sigma$-代数. 特别地, 包含 $\mathrm{F}_{\sigma}$ 集($\bigcup_{n=1}^{\infty} F_n$, $F_n$ 为闭集) 和 $\mathrm{G}_{\delta}$ 集($\bigcap_{n=1}^{\infty} G_n$, $G_n$ 为开集).
\end{example}
Borel 集是由 $\RR$ 中的开集生成的, 而对于 $\RR$ 中的开集$U$ 我们有如下引理
\begin{lemma}\label{lem:open-set-eq-open-inner-union}
    $U$ 是 $R$ 中的开集, 则存在至多可数个两两不交的开区间 $I_k$, 使得 $U = \bigcup_k I_k$.
\end{lemma}
\begin{proof}
    利用连通分支分解并在每个连通分支都能取一个有理数即可.
\end{proof}

\sout{\texorpdfstring{有人可能认为, Borel 集是一步到位的, 开集、闭集、可数并交搞几次就完了.}{}} \textbf{实际上}, Borel 集有一个严格的层级结构(Borel hierarchy): 
\begin{description}
    \item[第0层] 开集, 闭集
    \item[第1层] 可数交, 可数并, 如 $\mathrm{F}_{\sigma}$ 集, $G_\delta$ 集
    \item[第2层] $F_{\sigma\delta}$ ($F_\sigma$ 的可数交)、$G_{\delta\sigma}$($G_\delta$ 的可数并), 依此类推
    \item[极限层] 我们需要迭代到第一个\textbf{不可数序数($\omega_1$)} 才能把所有 Borel 集合穷尽。
\end{description}
这意味着存在一些 Borel 集合, 你需要操作不可数无穷多次（比如先并再交再并…重复 $\omega_1$ 次）才能得到。你无法用有限或可数次“开集、补集、可数并”的语言写出它的显式构造——虽然它是 Borel 的。

\begin{remark}
    如果取一个 Borel 集合 $B \subset \RR^2$，然后对它的投影（$\pi(B)=\{x:\exists y \st (x,y)\in B\}$），结果得到的集合不一定还是 Borel 集.
    可以参考\cite{CunZaiFeiBorelDeJieXiJiSuslinShiRuHeGeiLebesgueJiuCuoDe}.
\end{remark}
\begin{quote} \label{quote:borel-cardinality}
    尽管 Borel 集合看起来很庞大，但在基数意义上它非常少.
    具体来说, Borel 集合的基数不超过 $\mathfrak{c}$，而所有实数的集合的基数也是 $\mathfrak{c}$, 但是实数集的所有子集的个数是 $2^{\mathfrak{c}}$, 说明绝大多数实数集的子集都不是 Borel 的.
\end{quote}

接下来, 我们给一个Borel 集的经典例子:

\begin{example}[一个Borel 集的例子] \label{ex:contor-def}
    \textbf{康托(三分)集 (Cantor 集)}

    第一步, 把闭区间 $[0, 1]$ 三等分, 挖去中间的开区间 $(\frac{1}{3}, \frac{2}{3})$, 剩下的两个闭区间 $[0, \frac{1}{3}]$ 与 $[\frac{2}{3}, 1]$, 记这两个闭区间之并为 $F_1$. 
    第二步, 再把这两个闭区间三等分, 各挖去中间的两个开区间 $(\frac{1}{9}, \frac{2}{9})$ 与 $(\frac{7}{9}, \frac{8}{9})$, 记剩下的四个闭区间之并为 $F_2$.
    一般的, 当进行到第 $n$ 步时, 剩下 $2^n$ 个长度是 $3^{-n}$ 的互不相交的闭区间, 记这些闭区间之并为 $F_n$. 这样进行下去, 于是在 $[0, 1]$ 中挖去可数个互不相交的开区间, 而剩下来的点所成之集 $C$ 称为 Cantor 集.
    这时, 显然有 $C = \bigcap\limits_{n=1}^{\infty} F_n$.
\end{example}

\subsection{Cantor 集}
Cantor 集是一个非常经典的构造且有许多重要的性质需要了解.

\begin{proposition} \label{prop:cantor-closed}
    Cantor 集 $C$ 是非空有界闭集.
\end{proposition}
\begin{proof}
    因为每个 $F_n$ 是闭集以及闭集的任意交是闭集, 可知 $C = \bigcap\limits_{n=1}^\infty F_n$ 是闭集. 再因为 $\{F_n\}$ 是非空有界递减闭集列, 由闭集套定理知, $C = \bigcap\limits_{n=1}^\infty F_n \ne \emptyset$.
\end{proof}

\begin{proposition} \label{prop:cantor-perfect}
    $C = C'$.\quad 
\end{proposition}
\begin{proof}
设 $x \in C$, $\forall r>0$, 取 $n$ 使得 $3^{-n}<r$, 因为 $x\in F_n$, 而 $F_n$ 是一些长度为 $3^{-n}$ 的闭区间的并, 记其中含 $x$
的闭区间为 $[a, b]$, 显然有 $[a, b]\subset(x-r, x+r)$. 由 $C$ 的构造知, $a, b \in C$, 在 $a$ 与 $b$ 中至少有一个点不是 $x$,
所以 $((x-r, x+r)-\{x\})\cap C \ne \emptyset$, 即有 $x\in C'$, 所以 $C \subset C'$.

另一方面, 由\cref{prop:cantor-closed} $C$ 是闭集,  故有 $C \supset C'$. 综上, $C = C'$.
\end{proof}

对于像 Cantor 集这样, 导集为自身的集合称为完备集\footnote{完备集: perfect set, 有的地方称之为完美集.}.

\begin{proposition} \label{prop:cantor-has-no-inter-point}
    Cantor 集 $C$ 没有内点.
\end{proposition}
\begin{proof}
    对任意一开区间 $(a, b)$, 因为 $F_n$ 是一些长度为 $3^{-n}$ 的两两不交的闭区间的并, 当 $3^{-n} < b-a$ 时, 不能有 $(a, b) \subset F_n$,
    更不能有 $(a, b)\subset C$, 所以 $C$ 中不包含任何开区间, 即没有内点.
\end{proof}

\begin{proposition}
    Cantor 集 $C$ 具有连续基数 $\mathfrak{c}$.
\end{proposition}
\begin{proof}
    使用三进制小数表示 $[0, 1]$ 中的点时, Cantor 集中的点恰好不含数字 $1$, 即
    \[
    C = \{0.x_1x_2x_3\cdots x_k\cdots : x_k \in \{0, 2\}, k=1, 2, 3, \cdots\}.
    \]
    记 $E = \{(x_1, x_2, \dots, x_k, \dots) : x_k \in \{0, 1\}, k=1, 2, 3, \dots\}$ . 对 $x=0.x_1x_2x_3\cdots x_k\cdots \in C$, 令\[
    f(x)=(\frac{x_1}{2}, \frac{x_2}{2}, \cdots, \frac{x_k}{2}, \cdots).
    \]
    则 $f: C \to E$ 是一个双射. 从而知道 $C$ 的集合势等于 $E$ 的集合势 $\mathfrak{c}$.
\end{proof}
\begin{remark}
    值得注意的是, 在长度为 $1$ 的闭区间 $[0, 1]$ 中挖去可数个开区间, 这些开区间的长度和为 $3^{-1}+2\cdot 3^{-2} + \cdots + 2^n \cdots 3^{-n} + \cdots = 1$, 而剩下的 Cantor 集 $C$
    可以被看作 "长度" 为 $0$, 但是居然有连续基数 $\mathfrak{c}$.
\end{remark}

\subsection{稠密性}

接下来介绍一下稠密集和疏朗集\footnote{疏朗集: nowhere dense set, 有的地方称之为稀疏集或者无处稠密}的一些内容.

\begin{definition}
    若 $\overline{E} = \RR^n$, 称 $E$ 是 $\RR^n$ 中的稠密集. 若 $\left(\overline{E}\right)^{\circ} = \emptyset$, 称 $E$ 是 $\RR^n$ 中的疏朗集.
\end{definition}

\begin{lemma} \label{lem:dense-iff}
    $E$ 是 $\RR^n$ 中的稠密集 $\iff$ 对任意的非空开集 $G$ 有 $G \cap E \ne \emptyset$.
\end{lemma}

\begin{proposition}
    Cantor 集 $C$ 是 $\RR$ 中的疏朗集.
\end{proposition}
\begin{proof}
由 \cref{prop:cantor-closed} 和 \cref{prop:cantor-has-no-inter-point} 知, Cantor集是没有内点的闭集, 故 $\left(\overline E\right)^\circ = E^\circ = \emptyset$.
\end{proof}

\begin{lemma} 
    设 $G$ 是 $\RR^n$ 中的稠密开集, 则对于任意的非空开集 $S$, 存在开球 $B$ 使得 $\overline{B} \subset S \cap G$.
\end{lemma}
\begin{proof}
    对任意的非空开集 $S$, 有 \cref{lem:dense-iff} 知, $S\cap G$是非空开集. 取 $y\in S \cap G$, 存在 $r>0$ 使得
    $B(y, r) \subset S \cap G$. 令 $B = B(y, \frac{r}{2})$, 则 $\overline{B} \subset B(y, r) \subset S \cap G$.
\end{proof}

\begin{lemma}
    $\RR^n$ 中可数个稠密开集的交也是稠密集.
\end{lemma}
\begin{proof}
    设 $\{G_k\}$ 是 $\RR^n$ 中的稠密开集序列. 对任意的非空开集 $S$, 则由 $G_1$ 是稠密开集以及 \cref{lem:dense-iff} 知, 存在开球 $B_1$ 使得 
    $\overline{B_1}\subset S \cap G_1$. 再对开球 $B_1$ 以及稠密开集 $G_2$ 继续用引理,存在开球 $B_2$ 使得 $\overline{B_2} \subset B_1 \cap G_2$,
    继续下去得到闭球列:\[
    \overline{B_1} \supset \overline{B_2} \supset \overline{B_3} \supset \dots \supset \overline{B_k} \supset \dots
    \]
    由闭集套定理, 存在 $x \in \bigcap\limits_{k=1}^\infty \overline{B_k} \subset S$. 另一方面, 对每一个 $k$ 有 $x\in \overline{B_k} \subset G_k$, 故 $x \in \bigcap\limits_{k=1}^\infty G_k$,
    所以 $S \cap \left(\bigcap\limits_{k=1}^\infty G_k\right) \ne \emptyset$. 于是由 \cref{lem:dense-iff} 知$\bigcap\limits_{k=1}^\infty G_k$ 在 $\RR^n$ 中是稠密集.
\end{proof}
















\subsection{外测度}
设 $E \subset \RR$, 如何去度量 $E$ 的"长度"? 如果我们定义开区间 $(a, b)$ 的"长度" 是 $b-a$, 那么一个想法是使用开区间去覆盖和逼近$E$, 这也是外侧都的想法.
\begin{definition}[外测度]\label{def:exiter-meather}
    设 $E \subset \RR$, 定义 $m^*$ 为外测度, 
    $$m^*(A) = \inf\setof{\sum_{k=1}^{+\infty}l(I_k)}{\bigcup_{k=1}^{+\infty}I_k \supset A, I_k\text{为有界开区间}}$$
    其中若开区间 $I_k = (a_k, b_k)$, 定义 $l(I_k) = b_k - a_k$.
\end{definition}
事实上, 外测度就是一个从 $\RR$ 的幂集到 $\overline{\RR}_{\ge 0} = \RR_{\ge 0} \cup \{+\infty\}$ 的映射.
